% !TeX program = xelatex
% !TeX encoding = UTF-8
% 清华大学 IMDT 硕士学位论文开题报告
% 独立文件：在 Overleaf 中将本文件设为 Main document，选择 XeLaTeX。
% 依据用户提供的 IMDT Word 模板转换；非学校官方发布的 LaTeX 模板。
\documentclass[UTF8,fontset=fandol,zihao=-4,oneside,openany]{ctexbook}
\usepackage{iftex}
\ifXeTeX\else
  \PackageError{imdt}{Please select XeLaTeX as the compiler}{This template requires XeLaTeX.}
\fi
\usepackage[a4paper,left=31.75mm,right=31.75mm,top=25.4mm,bottom=25.4mm,
  headheight=15pt,headsep=7mm,footskip=10mm]{geometry}
\usepackage{fontspec}
\setmainfont{texgyretermes-regular.otf}[
  BoldFont=texgyretermes-bold.otf,
  ItalicFont=texgyretermes-italic.otf,
  BoldItalicFont=texgyretermes-bolditalic.otf]
\setsansfont{texgyreheros-regular.otf}[
  BoldFont=texgyreheros-bold.otf,
  ItalicFont=texgyreheros-italic.otf,
  BoldItalicFont=texgyreheros-bolditalic.otf]
% 默认字体随 TeX Live 提供，Overleaf 无需上传字体。
% 若院系要求指定的 Windows 字库，可合法上传对应字体后在此替换：
% \setCJKmainfont{SimSun} \setCJKsansfont{SimHei}
% \setmainfont{Times New Roman} \setsansfont{Arial}

\usepackage{amsmath,amssymb}
\usepackage{graphicx}
\usepackage{array,tabularx,booktabs}
\usepackage{caption}
\usepackage{enumitem}
\usepackage[super]{cite}
\usepackage{xcolor}
\usepackage{fancyhdr}
\usepackage{tocloft}
\usepackage[hidelinks,unicode,bookmarksnumbered]{hyperref}

% ==================== 个人信息填写区 ====================
\newcommand{\ThesisTitle}{请填写论文题目}
% 题目按原模板要求控制在 25 个汉字（符）以内。
\newcommand{\StudentName}{请填写姓名}
\newcommand{\StudentID}{请填写学号}
\newcommand{\SchoolName}{深圳国际研究生院}
\newcommand{\MajorName}{电子信息（互动媒体设计与技术）}
\newcommand{\Supervisor}{请填写导师姓名}
% 指导教师仅填写学籍系统中的导师。
\newcommand{\ReportDate}{\hspace{1em}年\hspace{1em}月\hspace{1em}日}
\newcommand{\ExpectedDefenseDate}{请填写预计答辩年月}

\newif\ifShowHints
\ShowHintstrue
% 完成填写后改成 \ShowHintsfalse，可隐藏所有 \hint{...} 写作提示。
\newcommand{\hint}[1]{\ifShowHints{\color{black!55}#1}\par\fi}

% ==================== 正文格式 ====================
% 原模板前部明确规定：正文 12bp、固定 20bp 行距；
% 后部范例曾写 1.5 倍行距，此处采用前部的详细规定。
\makeatletter
\renewcommand\normalsize{\@setfontsize\normalsize{12bp}{20bp}}
\makeatother
\AtBeginDocument{\normalsize}
\setlength{\parindent}{2em}
\setlength{\parskip}{0pt}
\linespread{1}
\setcounter{secnumdepth}{3}
\setcounter{tocdepth}{2} % 目录列至二级节标题，如 2.2.5
\ctexset{
  chapter = {
    name = {,}, number = \arabic{chapter},
    format = \centering\heiti\sffamily\fontsize{16bp}{20bp}\selectfont,
    beforeskip = 24bp, afterskip = 18bp, fixskip = true,
    aftername = \hspace{1em}, pagestyle = fancy
  },
  section = {
    format = \heiti\sffamily\fontsize{14bp}{20bp}\selectfont,
    beforeskip = 24bp, afterskip = 6bp, aftername = \hspace{1em}
  },
  subsection = {
    format = \heiti\sffamily\fontsize{13bp}{20bp}\selectfont,
    beforeskip = 12bp, afterskip = 6bp, aftername = \hspace{1em}
  },
  subsubsection = {
    format = \heiti\sffamily\fontsize{12bp}{20bp}\selectfont,
    beforeskip = 12bp, afterskip = 6bp, aftername = \hspace{1em}
  },
  contentsname = {目\quad 录},
  bibname = {参考文献}
}
\pagestyle{fancy}
\fancyhf{}
\fancyhead[C]{\songti\fontsize{10.5bp}{14bp}\selectfont\nouppercase{\leftmark}}
\fancyfoot[C]{\fontsize{10.5bp}{14bp}\selectfont\thepage}
\renewcommand{\headrulewidth}{0.4pt}
\renewcommand{\chaptermark}[1]{\markboth{\thechapter\quad #1}{}}
\fancypagestyle{plain}{%
  \fancyhf{}
  \fancyhead[C]{\songti\fontsize{10.5bp}{14bp}\selectfont\nouppercase{\leftmark}}
  \fancyfoot[C]{\fontsize{10.5bp}{14bp}\selectfont\thepage}
  \renewcommand{\headrulewidth}{0.4pt}
}

% 目录：章不缩进，一级节缩进 1 字，二级节缩进 2 字；20bp 行距。
\renewcommand{\cfttoctitlefont}{\hfill\heiti\sffamily\fontsize{16bp}{20bp}\selectfont}
\renewcommand{\cftaftertoctitle}{\hfill}
\setlength{\cftbeforetoctitleskip}{24bp}
\setlength{\cftaftertoctitleskip}{18bp}
\renewcommand{\cftchapfont}{\heiti\sffamily}
\renewcommand{\cftchappagefont}{\sffamily}
\renewcommand{\cftchapleader}{\cftdotfill{\cftdotsep}}
\setlength{\cftbeforechapskip}{0pt}
\setlength{\cftbeforesecskip}{0pt}
\setlength{\cftbeforesubsecskip}{0pt}
\cftsetindents{chapter}{0em}{2em}
\cftsetindents{section}{1em}{3em}
\cftsetindents{subsection}{2em}{4em}

% 图、表、公式按章编号：图 3.1、表 3.1、式（3.1）。
\renewcommand{\thefigure}{\thechapter.\arabic{figure}}
\renewcommand{\thetable}{\thechapter.\arabic{table}}
\renewcommand{\theequation}{\thechapter.\arabic{equation}}
\DeclareCaptionFont{imdt}{\songti\fontsize{11bp}{14bp}\selectfont}
\DeclareCaptionLabelSeparator{zhspace}{\hspace{1em}}
\captionsetup{font=imdt,labelsep=zhspace,justification=centering}
\captionsetup[figure]{position=bottom,aboveskip=6bp,belowskip=12bp}
\captionsetup[table]{position=top,aboveskip=6bp,belowskip=6bp}
\setlength{\heavyrulewidth}{1.5bp}
\setlength{\lightrulewidth}{1bp}
\newcolumntype{Y}{>{\centering\arraybackslash}X}
\setlist{nosep,leftmargin=2em}
\hypersetup{pdftitle={\ThesisTitle},pdfauthor={\StudentName}}

\begin{document}
\hypersetup{pageanchor=false}

% ==================== 第一页：封面 ====================
\begin{titlepage}
  \thispagestyle{empty}
  \centering
  \vspace*{14mm}
  {\heiti\fontsize{18bp}{26bp}\selectfont 清华大学硕士学位论文开题报告\par}
  \vspace{21mm}
  {\heiti\fontsize{22bp}{32bp}\selectfont \ThesisTitle\par}
  \vfill
  \begingroup
  \fontsize{14bp}{25bp}\selectfont
  \renewcommand{\arraystretch}{1.45}
  \begin{tabular}{@{}p{30mm}p{112mm}@{}}
    培养单位： & \SchoolName \\
    专\hspace{2em}业： & \MajorName \\
    申\hspace{0.5em}请\hspace{0.5em}人： & \StudentName \\
    学\hspace{2em}号： & \StudentID \\
    指导教师： & \Supervisor \\
  \end{tabular}
  \endgroup
  \vspace{25mm}

  {\fontsize{14bp}{20bp}\selectfont \ReportDate\par}
  \vspace{10mm}
\end{titlepage}
\hypersetup{pageanchor=true}

% ==================== 第二页：自动目录 ====================
\pagenumbering{Roman}
\pagestyle{plain}
\markboth{目录}{}
\tableofcontents
\clearpage
% ==================== 正文：从第 1 页开始编号 ====================
\pagenumbering{arabic}
\pagestyle{fancy}

% 原 Word 提供四章版与六章版大纲。本模板采用四章版，覆盖所有必需内容。
% 章标题用 \chapter，一级节用 \section，二级节用 \subsection。
% 标题自动编号；无需手动输入“第1章”或“1.1”。
\chapter{课题背景与意义}
\section{选题背景}
\hint{【在此说明研究所处的互动媒体设计与技术背景、具体应用场景、目标人群及现有问题。请用正式学术语言替换提示文字。】}

\section{研究意义}
\hint{【分别说明选题的理论意义、设计价值与工程应用价值，明确研究的边界与价值判断依据。】}

\chapter{国内外研究现状及分析}
\section{相关研究综述}
\hint{【按研究问题或技术路线梳理国内外研究，比较其方法、证据和局限。文献综述应进行归纳与分析，避免简单堆砌摘要。原模板要求查阅不少于20篇中外文献，书面报告不少于5000字。】}
\subsection{研究方向一}
\hint{【将本节标题替换为具体研究方向。用 \texttt{\string\cite\{reference-placeholder\}} 在正文中插入上标引用，如此处的占位引用\cite{reference-placeholder}；文末对应条目须替换为真实文献。】}
\subsection{研究方向二}
\hint{【总结另一相关研究方向，比较与本课题的关系。可按需要增加二级节标题。】}
\section{现有研究的不足与本课题切入点}
\hint{【从前文文献证据提出尚待解决的问题，并说明本课题如何回应这些问题。】}

\chapter{研究内容和主要创新点}
\section{研究目标与研究内容}
\hint{【明确总体目标，提炼3至4个可实施、可验证的研究点，并说明各研究点之间的关系。】}

\section{研究方案与研究方法}
\subsection{总体方案与技术路线}
\hint{【说明设计流程或技术路线。下图为占位示例，引用外部图片须注明出处。】}
\begin{figure}[htbp]
  \centering
  % 有图片后，用下行替换占位框，并将图片上传到 Overleaf：
  % \includegraphics[width=0.85\linewidth]{your-figure.png}
  \fbox{\parbox[c][20mm][c]{0.85\linewidth}{\centering
    在此放置研究技术路线或系统架构图\\
    图中文字、术语与单位应与正文保持一致}}
  \caption{研究技术路线示意图（待替换）}
  \label{fig:research-route}
\end{figure}

\subsection{研究方法与评价方案}
\hint{【说明研究方法、数据来源、实验或用户研究流程、评价指标及可行性。式（\ref{eq:example}）为编号示例，提交前替换或删除。】}
\begin{equation}
  \label{eq:example}
  \bar{x}=\frac{1}{n}\sum_{i=1}^{n}x_i
\end{equation}

\section{主要创新点}
\hint{【对照已有研究，说明预期的设计或技术贡献及其验证依据。】}

\section{关键难点及拟采取的解决措施}
\hint{【逐项说明关键难点、产生原因、应对措施、替代方案及验证方法。】}

\section{已具备和所需的条件}
\hint{【说明已有工作基础、阶段性进展、设备和实验条件、数据与软件资源、外协计划及经费；明确尚待解决的条件问题。】}

\chapter{预期成果及论文进度安排}
\section{预期成果}
\hint{【说明预期达到的水平和成果形式，如论文、交互原型、软件系统或设计作品，并给出完成或验收依据。】}

\section{论文进度安排}
\hint{【将表中的占位内容替换为具体年月和可检查的阶段成果。表题在表格上方，图题在图片下方；三线表与编号已设置。】}
\begin{table}[htbp]
  \caption{论文工作进度安排}
  \label{tab:schedule}
  \centering
  \fontsize{11bp}{14bp}\selectfont
  \renewcommand{\arraystretch}{1.45}
  \begin{tabularx}{\linewidth}{@{}p{24mm}XY@{}}
    \toprule
    时间阶段 & 主要工作 & 阶段成果 \\
    \midrule
    待填起止年月 & 文献阅读与需求调研 & 文献综述与问题定义 \\
    待填起止年月 & 方案设计与原型实现 & 设计方案与可运行原型 \\
    待填起止年月 & 实验验证与设计迭代 & 分析结果与迭代记录 \\
    待填起止年月 & 论文撰写与修改 & 学位论文与成果材料 \\
    \bottomrule
  \end{tabularx}
\end{table}

\section{预计答辩时间}
预计答辩时间：\ExpectedDefenseDate。

% ==================== 参考文献 ====================
% 为保持两个 .tex 均为独立文件，此处使用内置 thebibliography，
% 不依赖额外 .bib、.bst、.cls 文件，也无需运行 Biber。
% 按正文首次引用顺序排列条目，并按最新写作指南人工核对著录格式。
% 本示例不是实际文献，必须替换；继续添加 \bibitem{唯一键} 即可。
\begin{thebibliography}{99}
  \addcontentsline{toc}{chapter}{参考文献}
  \markboth{参考文献}{}
  \fontsize{10.5bp}{16bp}\selectfont
  \setlength{\itemsep}{3bp}
  \setlength{\parskip}{0pt}
  \bibitem{reference-placeholder}
  【文献条目占位，提交前替换】作者. 文献题名[J]. 期刊名, 年, 卷(期): 起止页码.
\end{thebibliography}

\end{document}
